\section{Follow-up: are the two edits specific?}
\label{sec:followup}

% Design text: proposal rev 1b, manifest b2b14bac, committed a177536 before held output.
% Leave results pending until the Advisor recount is complete.

The confirmation leaves open whether the selected components and directions matter, or whether other edits
of the same kind would work as well. A follow-up was pre-registered after confirmation closed. It uses a new
controlled split, N (seed 54), disjoint from the agent's exploration and the selection split. The follow-up
addresses three questions without changing the confirmation verdicts.

\paragraph{Does the read direction matter?}
The agent described the layer-35 edit as reading an ``answer available'' signal and writing along
city~$-$~``unknown''. Twenty control edits keep the write vector but use a random read vector. Each is scaled
on development rows to match the edit's root-mean-square gate and sign-aligned to push in the same direction.
The frozen criterion, C4, requires the edit's utility to exceed all twenty controls and its unsupported-answer
rate to fall below their median. Failure would leave read-side specificity unestablished; it would not by
itself show equivalence to an output bias.

\paragraph{Are the two heads unusual among late-layer heads?}
C1 requires the target pair's ablation to recover more false abstentions than each of twenty random pairs
from the same layers, with a margin of at least 0.30 over their median. C2 requires doubling the target
pair to create more abstentions on previously correct rows than doubling each random pair.

\paragraph{Does the layer-35 signal follow the agent's predictions?}
C5 tests four frozen predictions about changes in the proposed signal when evidence is removed, the city
changes, distractors are added, or unrelated evidence replaces absent evidence.

C1, C2 and C4 use an empirical rank of 1 among 21 evaluated edits; C5 uses an intersection--union test at
0.05. No family-wise error-control claim is made. Under a null model in which target scores are exchangeable
with the random-control scores, with ties counted against the target, and with a valid C5 test, a union bound
puts the chance of at least one passing claim at most about 0.19. This error bound applies under the stated
exchangeability null and a valid C5 test; it is not a guarantee for every possible mechanism of false abstention.

Descriptive arms examine single heads, other late-layer pairs, random write vectors, direct logit attribution,
and the prompt with its ``unknown'' option removed. A natural-question check uses 138 LongMemEval development
questions with evidence in view.

\paragraph{Results.} \pending{§53 results: C1, C2, C4, C5 verdicts; LongMemEval releases, damage and answers on
abstention questions; spend.}
